% Generated draft: requires mathematical and source review.
\documentclass[11pt]{article}
\usepackage[margin=25mm]{geometry}
\usepackage{fontspec}
\setmainfont{Latin Modern Roman}
\usepackage{amsmath,amssymb,mathtools,unicode-math}
\setmathfont{Latin Modern Math}
\usepackage{enumitem,needspace,xurl,hyperref}
\setlength{\parindent}{0pt}
\setlength{\parskip}{5pt}
\setlength{\emergencystretch}{3em}
\newcommand{\N}{\mathbb{N}}
\newcommand{\Z}{\mathbb{Z}}
\newcommand{\Q}{\mathbb{Q}}
\newcommand{\R}{\mathbb{R}}
\newcommand{\C}{\mathbb{C}}
\newcommand{\F}{\mathbb{F}}
\newcommand{\A}{\mathbb{A}}
\newcommand{\PP}{\mathbb{P}}
\newcommand{\E}{\mathbb{E}}
\newcommand{\eps}{\varepsilon}
\newcommand{\dd}{\,\mathrm{d}}
\begin{document}
\section*{UnsolvedMath ID: 30000137}
\textbf{Source problem code:} catalogue entry
\subsection{Definitions and mathematical statement}
For a \(C^1\) embedding \(f:M^d\hookrightarrow\R^N\), let \(E_p=df_p(T_pM)\). The embedding is totally skew if, for every pair of distinct points \(p,q\in M\), every tangent line at \(f(p)\) and every tangent line at \(f(q)\) are disjoint and nonparallel. Equivalently,
\[
\dim\bigl(E_p+E_q+\operatorname{span}\{f(q)-f(p)\}\bigr)=2d+1.
\]
Thus the two tangent direction spaces and the connecting direction must be linearly independent as subspaces; disjointness of tangent spaces alone is not the full condition. Define
\[
N_{\mathrm{skew}}(M)=\min\{N:\text{there is a totally skew embedding }M\hookrightarrow\R^N\},
\]
with value infinity if the set is empty. The general source question asks whether the cited lower bounds for \(N_{\mathrm{skew}}(M)\) are sharp. No particular lower-bound formula or range of manifolds is reproduced, so this part remains a source-dependent sharpness programme.

The explicit test case is whether there exists a totally skew embedding
\[
f:D^3\hookrightarrow\R^7.
\]
Using the closed-disk convention \(D^3=\{x\in\R^3:\|x\|\leq1\}\), this asks for an embedding differentiable up to the boundary such that for every distinct \(p,q\in D^3\),
\[
f(q)-f(p),\ \partial_1f(p),\ \partial_2f(p),\ \partial_3f(p),\
\partial_1f(q),\ \partial_2f(q),\ \partial_3f(q)
\]
are seven linearly independent vectors in \(\R^7\). Boundary tangent spaces have dimension three. The source does not specify the regularity class or whether disk denotes a closed disk, open disk, or local embedded disk; these conventions require confirmation, and the displayed test case makes its convention explicit.

The supplied literature assessment checked 2026-08-21 says that the totally skew three-disc in \(\R^7\) remains listed as an unresolved case in later treatments of the Ghomi--Tabachnikov bounds.
\subsection{Short English statement}
Ask whether known dimension lower bounds for totally skew embeddings are attained, particularly whether a three-dimensional disk embeds in \(\R^7\) with tangent lines at distinct points always disjoint and nonparallel.
\subsection{Sources}
\begin{itemize}
\item Local UnsolvedMath registry, permanent ID 30000137, catalogue entry, \emph{Sharp Dimensions for Totally Skew Embeddings}.
\item \emph{Geometrie} (2005), Oberwolfach Reports, pp. 2491--2538, DOI \url{https://doi.org/10.4171/owr/2004/47}.
\item Supplied reference labeled ``totally skew embeddings paper'': \url{https://ghomi.math.gatech.edu/Papers/totallyskew.pdf}; the assessment associates the bounds with Ghomi--Tabachnikov.
\item Supplied statement verification reports confirmation of both the sharpness question and three-disk case; supplied literature triage is dated 2026-08-21.
\end{itemize}
\end{document}
